\origpage{106}{ROBERTSON AND SEYMOUR}
\noindent be the walk $v_{i_1}, e_{i_1 + 1}, \dots, v_{i_2}$ and let $W_2$ be the walk $v_{i_3}, e_{i_3 + 1}, \dots, v_{i_4}$. Then no internal vertex of $W_1$ or $W_2$ lies in $A \cup B$, and no edge of $W_1$ or $W_2$ lies in $A$ or in $B$. Hence the sum of the lengths of $W_1$ and $W_2$ is at most $l(Y) - \alpha - \beta$, where $\alpha$ is the number of values of $i(1 \le i \le t)$ with $e_i \subseteq \text{bd}(A)$, and $\beta$ is defined similarly. We may assume then that $W_1$ has length $\le \frac{1}{2}(l(Y) - \alpha - \beta)$. Now $v_{i_1} \ne v_{i_2}$ since $A \cap B = \emptyset$; choose a line $F$ with ends $v_{i_1}, v_{i_2}$ and with $F \subseteq U(W_1)$. Then $X = A \cup B \cup F$ is the required tie; for if $\gamma$ is the number of edges of $\Gamma$ included in $F$, then $\gamma \le \frac{1}{2}(l(Y) - \alpha - \beta)$ and $l(X) = \alpha + \beta + 2\gamma$, and so $l(X) \le l(Y)$. $\square$

\sectionhead{References}

1. D. B. A. Epstein, Curves on 2-manifolds and isotopies, \textit{Acta Math.} \textbf{115} (1966), 83-107. 2. N. Robertson and P. D. Seymour, Graph minors. III. Planar tree-width, \textit{J. Combin. Theory Ser. B} \textbf{36} (1984), 49-64. 3. N. Robertson and P. D. Seymour, Graph minors. VII. Disjoint paths on a surface, \textit{J. Combin. Theory Ser. B} \textbf{45} (1988), 212-254. 4. N. Robertson and P. D. Seymour, Graph minors. X. Obstructions to tree-decomposition, \textit{J. Combin. Theory Ser. B} \textbf{52} (1991), 153-190. 5. P. D. Seymour and R. Thomas, Call routing and the ratcatcher, \textit{Combinatorica}, to appear.
